\documentclass[11pt,a4paper]{article}
\usepackage[margin=1in]{geometry}
\usepackage{kotex}
\usepackage{amsmath,amssymb,amsfonts,bm,mathtools}
\usepackage{booktabs,multirow,array,longtable}
\usepackage{enumitem}
\usepackage{graphicx}
\usepackage{hyperref}
\usepackage{natbib}
\usepackage{xcolor}
\usepackage{setspace}
\usepackage{caption}
\usepackage{float}
\usepackage{algorithm}
\usepackage{algpseudocode}
\usepackage{tcolorbox}
\hypersetup{colorlinks=true, linkcolor=blue, citecolor=blue, urlcolor=blue}
\setstretch{1.15}

\newtcolorbox{keybox}[1][]{colback=blue!5!white,colframe=blue!60!black,
  fonttitle=\bfseries,title={#1}}
\newtcolorbox{warnbox}[1][]{colback=red!5!white,colframe=red!60!black,
  fonttitle=\bfseries,title={#1}}

\title{PV 예측 연구 계획 v2:\\
시나리오 기반 베이지안 계층모형과 확률론적 화력 백업 최적화 연계}
\author{한국남동발전 태양광 변동성 대비 플랫폼}
\date{2026년 4월}

\begin{document}
\maketitle

\begin{abstract}
본 계획서는 \texttt{plan/pv/v1}에서 제시한 베이지안 계층모형 + FiLM-조건부 GRU decoder-only 프레임워크를 \textbf{시나리오 생성기(scenario generator)} 관점에서 재설계한 v2안이다. v1은 베이지안 1단계를 ``점예측 baseline 추출기''로 사용하였으나, 본 플랫폼의 궁극 목적이 \textit{태양광 변동성 대비 화력 백업 최적화}임을 고려하면 단일 점예측보다 \textbf{미래의 확률분포 자체}를 하류 최적화 모듈(Stochastic Unit Commitment, Chance-Constrained UC)에 전달하는 것이 결정적으로 중요하다. 따라서 v2는 (1) posterior predictive distribution을 전면에 내세우고, (2) AR(1) 관측 노이즈, (3) heteroscedastic precision, (4) Gaussian Process 공간 prior, (5) Zero-Inflated Beta (ZIB)를 도입하여 모델이 시간/공간/날씨 의존적 불확실성을 스스로 인식하도록 강화한다. 최종적으로 1,000개 이상의 완전한 포트폴리오 시계열 시나리오가 산출되며, 이 시나리오 번들이 D-1일 KPX 가용용량 신고 및 화력 예비력 산정의 직접적 입력이 된다.
\end{abstract}

\tableofcontents
\newpage

%=========================================================
\section{v1에서 v2로: 연구 철학의 전환}
%=========================================================

\subsection{왜 전환하는가}

v1은 논문 \texttt{pv\_predict.md}의 2단계 프레임워크를 충실히 구현하는 데 초점이 있었다. 그러나 본 플랫폼 운영자가 실제 D-1일 10시에 수행하는 의사결정은 ``내일의 태양광 발전량을 정확히 맞히는 것''이 아니라, \textbf{``어떤 화력 조합으로 내일의 예측 불확실성을 커버할 것인가''}이다. 즉, 플랫폼의 가치는 점예측 정확도가 아니라 \textbf{불확실성의 정량화 품질}에서 나온다.

\begin{keybox}[핵심 재정의]
\textit{우리가 모르는 부분을 확률로 표현하고, 그 확률분포에서 직접 추출한 시나리오를 화력 최적화에 공급한다.} 이 관점에서 PV 예측 모듈의 출력은 ``발전량 한 숫자''가 아니라 \textbf{``그럴듯한 내일 1,000가지''}이다.
\end{keybox}

\subsection{v1 vs v2 비교}

\begin{table}[H]
\centering
\small
\caption{v1과 v2의 설계 축 비교}
\begin{tabular}{lll}
\toprule
설계 축 & v1 (논문 충실 구현) & v2 (시나리오 중심 강화) \\
\midrule
1단계 목적 & 점예측 baseline $\hat\mu^{\mathrm{Bayes}}$ 추출 & 시나리오 번들 $\{\tilde y^{(s)}\}_{s=1}^{S}$ 생성 \\
관측 모형 & Beta (고정 $\phi_r$) & ZIB + heteroscedastic $\phi_{r,t}$ \\
시간 구조 & Fourier basis만 & Fourier + AR(1) 관측 노이즈 \\
공간 구조 & 지역별 random effect & 지역 + 사이트간 GP spatial prior \\
2단계 출력 & 잔차 점예측 & 잔차 + 분산 (probabilistic head) \\
평가 지표 & nRMSE, nMAE & + CRPS, coverage, pinball, ramp-hit, UC-cost \\
하류 연계 & 명시 안 됨 & Stochastic UC / Chance-Constrained UC 직결 \\
출력 차원 & $(H, N)$ 스칼라 텐서 & $(S, H, N)$ 시나리오 텐서 + 분위수 요약 \\
\bottomrule
\end{tabular}
\end{table}

\subsection{v2가 바꾸지 않는 것}

다음은 v1과 동일하게 유지된다.
\begin{itemize}[leftmargin=2em]
    \item 2단계 프레임 자체 (베이지안 1단계 + FiLM-GRU 2단계 직렬구조)
    \item 지역별 부분풀링(partial pooling)과 regional statistical signature 개념
    \item Fourier basis 기반 시간/계절 구조 모델링
    \item 논문 \texttt{pv\_predict.md}의 표기법과 식 번호 체계
\end{itemize}

즉, v2는 v1의 \textbf{대체가 아니라 확장}이며, v1의 수식은 v2에 특수 케이스로 포함된다.

%=========================================================
\section{1층과 2층의 역할 분담 (명시화)}
\label{sec:role_split}
%=========================================================

본 프레임워크의 본질은 \textbf{``1층이 통계적 구조를 짜내고, 2층이 비선형 잔차를 보정한다''}는 분업이다. 이 절은 이 분업을 한 표로 명시한다. v2의 강화 1--4가 ``1층이 2층의 일을 가져가는 것''이 아니라, ``1층의 구조 부분을 더 정교하게 만들어 2층이 진짜 비선형 잔차에 집중할 수 있게 한다''는 점을 명확히 한다.

\subsection{분업 원칙}

\begin{equation*}
    \underbrace{y_{i,t}}_{\text{관측}} = \underbrace{\mu^{\mathrm{Bayes}}_{i,t}}_{\text{1층: 구조}} + \underbrace{e_{i,t}}_{\text{잔차}},\qquad
    \underbrace{\hat e_{i,t+h}}_{\text{2층: 비선형 보정}} = f_{\mathrm{GRU+FiLM}}(\text{과거 } e, \text{기상}, \text{미래 baseline}).
\end{equation*}

\subsection{역할 분담 표}

\begin{table}[H]
\centering
\small
\caption{1층(베이지안)과 2층(GRU+FiLM)의 책임 분담}
\begin{tabular}{p{4.5cm}cc p{5.5cm}}
\toprule
시계열 성분 & 1층 & 2층 & 책임 분배 근거 \\
\midrule
지역별/사이트별 평균 레벨 & \checkmark & & 부분풀링으로 안정 추정 \\
일중/연중 주기 (Fourier) & \checkmark & & 명시적 주기함수가 더 적합 \\
기상 변수에 대한 \emph{선형} 반응 & \checkmark & & random slope $b_{r,\cdot}$로 처리 \\
시간 자기상관 (평균적 강도) & \checkmark & & AR(1) latent으로 시나리오 생성 시 필수 \\
사이트간 공간 상관 (평균적 강도) & \checkmark & & GP prior로 포트폴리오 분산 정확화 \\
날씨별 분산 차이 (이분산) & \checkmark & & heteroscedastic $\phi_{r,t}$ \\
일출/일몰 경계 행동 (질량 0) & \checkmark & & ZIB 관측분포 \\
\midrule
기상-PV \emph{비선형} 상호작용 & & \checkmark & 1층은 logit-선형, GRU가 비선형 \\
단기 급변(ramp) 미세 패턴 & & \checkmark & 시점간 비선형 의존성 \\
과거 잔차의 비선형 시차 효과 & & \checkmark & GRU의 hidden state가 요약 \\
미래 외생변수와 과거의 비선형 결합 & & \checkmark & decoder-only head가 결합 \\
사이트별 미세기상 잔여 패턴 & & \checkmark & FiLM이 지역 조건부로 보정 \\
\midrule
파라미터 불확실성 (epistemic) 정량화 & \checkmark & & posterior 자체가 분포 \\
관측 노이즈 (aleatoric) 시나리오화 & \checkmark & & ZIB + AR(1) 샘플링 \\
2층 잔차의 분산 정량화 & & \checkmark & probabilistic head (mean+var) \\
\bottomrule
\end{tabular}
\end{table}

\subsection{왜 시간 상관/공간 상관/이분산이 1층 책임인가}

이 점이 가장 헷갈리는 부분이므로 명시한다. 강화 1--3은 모두 \textbf{``평균적 통계 구조''}이며, 이는 \emph{시나리오 생성기}가 그럴듯한 시계열을 산출하기 위해 1층에서 필수적이다.

\begin{itemize}[leftmargin=2em]
    \item AR(1)이 1층에 없으면 → 한 시나리오 안의 시점들이 독립 → 비현실적 ``깜빡이는'' 시계열만 나옴 → 램프 시나리오 생성 불가
    \item GP가 1층에 없으면 → 사이트들이 독립 → 포트폴리오 합의 분산 과소평가 → 화력 예비력 부족
    \item 이분산이 1층에 없으면 → 흐린 날도 맑은 날도 같은 분산 → 분위수 캘리브레이션 실패
\end{itemize}

반면 2층은 \textbf{개별 점예측에서 비선형 잔차를 깎는 일}을 한다. 즉 ``평균이 아닌 곳, 비선형이 있는 곳''이 2층의 영역이다.

\begin{keybox}[한 줄로]
1층 = ``평균적으로 어떻게 움직이는가''의 통계적 구조 모형 + 시나리오 생성기 \\
2층 = ``이 순간 비선형으로 어디로 튀는가''의 잔차 보정기
\end{keybox}

%=========================================================
\section{현재 보유 데이터와 가용성 매트릭스}
\label{sec:data_avail}
%=========================================================

\subsection{보유 데이터 인벤토리 (2026년 4월 기준)}

\begin{table}[H]
\centering
\small
\caption{현재 보유 데이터}
\begin{tabular}{p{3cm}p{3cm}p{3cm}p{4cm}}
\toprule
데이터 & 기간 & 해상도 & 비고 \\
\midrule
태양광 발전실적 & 2022-01 \textasciitilde{} 2025-03 & 사이트 \texttimes{} 호기 \texttimes{} 시간 & 39개월, 16개 사이트 \\
ASOS 기상 관측 & 2022-01 \textasciitilde{} 2025-03 & 관측소 \texttimes{} 시간 & 일사량(icsr), 전운량(dc10Tca), 기온, 습도, 풍속 모두 있음 \\
화력 발전실적 & 2022-01 \textasciitilde{} 2025-03 & 호기 \texttimes{} 시간 & 화력 단계용 \\
사이트 메타데이터 & 정적 & 사이트별 & 용량, 주소 있음. \textbf{위경도 비어있음} \\
사이트 현장 기상 & 2024-07 (1개월) & 사업소 \texttimes{} 일 & \textbf{시간대별 아님, 일사량 없음} → 사용 곤란 \\
LDAPS 예보 & 미수집 & --- & 활용신청 필요 \\
SMP, 입찰정보 & 보유 & --- & 본 모델 입력 아님 \\
\bottomrule
\end{tabular}
\end{table}

\subsection{데이터-모델 가용성 매트릭스}

각 강화 모듈이 현재 데이터로 학습 가능한지 진단한다.

\begin{table}[H]
\centering
\small
\caption{모듈별 데이터 요건과 현재 가용성}
\begin{tabular}{p{3.5cm}p{4cm}cp{4.5cm}}
\toprule
모듈 (1층 구성요소) & 필요 데이터 & 가용 & 보완 작업 \\
\midrule
Beta 회귀 baseline & 사이트별 시간 PV + GHI/T/CLD & \checkmark & ASOS와 사이트 매칭 (가까운 관측소) \\
지역별 random effect & 사이트--지역 매핑 & \checkmark & 메타에 region 컬럼 있음 \\
사이트별 random intercept & 사이트 ID & \checkmark & --- \\
일중/연중 Fourier & 시간/날짜 & \checkmark & 자동 계산 \\
Clear-sky GHI (CSI 계산) & 위경도 + 시각 & $\triangle$ & \textbf{geocoding 필요} (주소만 있음) \\
SZA (주야 분리, ZIB의 $\pi_0$) & 위경도 + 시각 & $\triangle$ & 동일 (geocoding 필요) \\
\midrule
\textbf{강화 1: AR(1)} & 시간 연속 PV 시계열 & \checkmark & 추가 작업 없음 \\
\textbf{강화 2: Heteroscedastic $\phi$} & GHI/CLD/CSI 시간 시계열 & \checkmark & ASOS의 icsr/dc10Tca로 충분 \\
\textbf{강화 3: GP 공간 prior} & 사이트별 정확한 위경도 & \textbf{\texttimes} & \textbf{\,geocoding 필수 (Naver/Kakao API)} \\
\textbf{강화 4: ZIB} & SZA 계산을 위한 위경도 & $\triangle$ & geocoding 후 가능 \\
\midrule
2층 GRU + FiLM & 1층 출력 + 과거/미래 기상 & \checkmark & 1층 완료 후 자연 가능 \\
실시간 보정 & LDAPS 또는 초단기예보 & \textbf{\texttimes} & \textbf{LDAPS 활용신청 필요} \\
\bottomrule
\end{tabular}
\end{table}

\subsection{즉시 해야 할 데이터 작업}

위 매트릭스에서 \textbf{\texttimes} 표시된 두 항목이 v2 구현의 전제조건이다.

\begin{enumerate}[leftmargin=2em]
    \item \textbf{사이트 16개 위경도 확보}: \texttt{data/solar\_sites.csv}의 주소 컬럼을 입력으로 Kakao Local API 또는 Naver Geocoding API로 일괄 변환. 작업량 1시간 이내. 이후 \texttt{lat, lon} 컬럼 채워넣기.
    \item \textbf{사이트 \texttimes{} ASOS 관측소 매핑 테이블}: 각 사이트에 대해 가장 가까운 ASOS 관측소를 1--3개 선정. 영흥 → 인천(112), 삼천포 → 진주(192) 또는 통영(162), 예천 → 안동(136) 등. 거리 가중 평균으로 사이트 기상 대용 변수 생성.
    \item \textbf{시간대별 결측 사이트 제외 또는 별도 처리}: \texttt{탑선태양광}, \texttt{여수태양광} 등은 시간대별 0이고 일별 합계만 들어와 1층 학습 대상에서 제외. \textbf{학습 대상 사이트 약 12--13개}.
    \item \textbf{LDAPS 활용신청}: 학습 자체는 ASOS 실측으로 가능하나, 실시간 D-1 운영 시 미래 기상 입력이 필요하므로 본 신청은 병렬로 진행.
\end{enumerate}

\subsection{데이터 한계로 인한 v2 모델 적용 우선순위 조정}

위 가용성에 근거하여 \textbf{구현 순서를 다음과 같이 조정}한다.

\begin{table}[H]
\centering
\small
\caption{데이터 가용성 기반 구현 우선순위 (재조정)}
\begin{tabular}{cp{6cm}p{5cm}}
\toprule
순위 & 작업 & 즉시 가능 여부 \\
\midrule
1 & geocoding으로 위경도 확보 & \checkmark 즉시 가능 \\
2 & ASOS-사이트 매핑 + 사이트별 기상 시계열 생성 & \checkmark 즉시 가능 \\
3 & v1 기준 1층 (Beta + random effect + Fourier) PyMC 구현 & \checkmark 즉시 가능 \\
4 & 강화 2 (heteroscedastic $\phi$) 추가 & \checkmark 즉시 가능 \\
5 & 강화 1 (AR(1)) 추가 & \checkmark 즉시 가능 \\
6 & 강화 4 (ZIB) 추가 & \checkmark geocoding 후 \\
7 & 강화 3 (GP 공간 prior) 추가 & \checkmark geocoding 후 \\
8 & 시나리오 번들 생성 + 분위수 요약 산출 & \checkmark 1층 완료 후 \\
9 & 2층 FiLM-GRU + probabilistic head & \checkmark 1층 완료 후 \\
10 & 실시간 운영 모드 (LDAPS) & \texttimes LDAPS 활용신청 후 \\
\bottomrule
\end{tabular}
\end{table}

\textbf{현 시점에서 LDAPS만 빼고 모두 즉시 착수 가능}하다. v1의 ``기상청 API 활용신청 대기'' 가정과 달리, ASOS 실측만으로도 학습 데이터셋을 구성할 수 있다는 점이 본 가용성 검토의 가장 중요한 발견이다.

%=========================================================
\section{수학적 기초: Posterior Predictive Distribution의 전면화}
%=========================================================

\subsection{재정의된 1단계 출력}

v1의 1단계는 posterior predictive mean만 사용한다.
\begin{equation}
    \hat\mu^{\mathrm{Bayes}}_{i,t+h} = \mathbb{E}[\tilde y_{i,t+h}\mid\mathcal{D}_{\mathrm{train}}]
    \approx \frac{1}{S}\sum_{s=1}^{S}\mu^{(s)}_{i,t+h}.
\end{equation}

v2의 1단계는 \textbf{full posterior predictive distribution}을 산출한다.
\begin{equation}
    p(\tilde y_{i,t+h}\mid\mathcal{D}_{\mathrm{train}}) = \int \underbrace{p(\tilde y_{i,t+h}\mid\theta)}_{\substack{\text{관측 노이즈}\\(\text{aleatoric})}}\,\underbrace{p(\theta\mid\mathcal{D}_{\mathrm{train}})}_{\substack{\text{파라미터 불확실성}\\(\text{epistemic})}}\,d\theta.
    \label{eq:posterior_predictive}
\end{equation}

MCMC 샘플 $\{\theta^{(s)}\}_{s=1}^{S}$로부터 식 \eqref{eq:posterior_predictive}의 샘플을 얻는 절차는 다음과 같다.

\begin{algorithm}[H]
\caption{Posterior Predictive Scenario Sampling (AR(1) latent state 명시)}
\label{alg:scenario}
\begin{algorithmic}[1]
\State MCMC 실행 후 posterior 샘플 $\{\theta^{(s)}\}_{s=1}^{S}$ 확보
\State $\theta = \{\alpha_0, a_r, u_i, b_{r,\cdot}, \beta_\cdot, \gamma_\cdot, \delta_\cdot, \rho_r, \sigma_{\nu,r}, \zeta_{r,\cdot}, \omega_{r,\cdot}\}$
\For{$s=1,\dots,S$}
    \State $\epsilon^{(s)}_{i,t} \gets$ 학습 구간 종점의 latent residual (Kalman/forward-filter로 추정 또는 0)
    \For{$h=1,\dots,H$}
        \State \emph{[1] 결정론적 부분 평가}
        \State $\eta^{(s)}_{i,t+h} \gets$ 식 \eqref{eq:bayes_stage1_v2} 평가 (파라미터 $\theta^{(s)}$, $\epsilon$ 제외)
        \State $\log\phi^{(s)}_{r,t+h} \gets \zeta^{(s)}_{r,0} + \zeta^{(s)}_{r,\mathrm{CLD}}\widetilde{\mathrm{CLD}}_{i,t+h} + \cdots$
        \State $\mathrm{logit}(\pi^{(s)}_{0,i,t+h}) \gets \omega^{(s)}_{r,0} + \omega^{(s)}_{r,\mathrm{SZA}}\,\mathrm{SZA}_{t+h}$
        \State \emph{[2] AR(1) latent 전진 샘플링}
        \State $\nu^{(s)}_{i,t+h} \sim \mathcal{N}(0,\,(\sigma^{(s)}_{\nu,r})^2)$
        \State $\epsilon^{(s)}_{i,t+h} \gets \rho^{(s)}_r \cdot \epsilon^{(s)}_{i,t+h-1} + \nu^{(s)}_{i,t+h}$
        \State \emph{[3] 평균 결합 후 ZIB 샘플링}
        \State $\mu^{(s)}_{i,t+h} \gets \sigma\bigl(\eta^{(s)}_{i,t+h} + \epsilon^{(s)}_{i,t+h}\bigr)$
        \State $z^{(s)}_{i,t+h} \sim \mathrm{Bernoulli}(\pi^{(s)}_{0,i,t+h})$
        \If{$z^{(s)}_{i,t+h}=1$}
            \State $\tilde y^{(s)}_{i,t+h} \gets 0$
        \Else
            \State $\tilde y^{(s)}_{i,t+h} \sim \mathrm{Beta}\bigl(\mu^{(s)}_{i,t+h}\phi^{(s)}_{r,t+h},\,(1-\mu^{(s)}_{i,t+h})\phi^{(s)}_{r,t+h}\bigr)$
        \EndIf
    \EndFor
\EndFor
\State \Return $\mathcal{S} = \{\tilde y^{(s)}_{i,t+1:t+H}\}_{s=1}^{S}$ \Comment{shape: $S\times H\times N$}
\end{algorithmic}
\end{algorithm}

핵심은 \textbf{2중 샘플링}이다. posterior $\theta^{(s)}$에서 한 번 뽑고(파라미터 불확실성), 그 $\theta^{(s)}$ 조건 하의 관측 분포에서 또 한 번 뽑는다(관측 노이즈). 또한 step [2]의 AR(1) 전진 샘플링이 \emph{같은 $s$ 인덱스 내에서} 시간을 가로지르며 진행되므로, 한 시나리오 안의 $H$개 시점이 자동으로 시간 상관을 갖는다. 결과적으로 식 \eqref{eq:posterior_predictive}의 샘플이 얻어진다.

\subsection{불확실성 분해}

모델은 다음 두 종류의 불확실성을 명시적으로 분리 보고한다.

\begin{align}
    \sigma^2_{\mathrm{total}}(t+h) &= \underbrace{\mathrm{Var}_{\theta}\bigl[\mathbb{E}[\tilde y_{t+h}\mid\theta]\bigr]}_{\sigma^2_{\mathrm{epistemic}}} + \underbrace{\mathbb{E}_{\theta}\bigl[\mathrm{Var}[\tilde y_{t+h}\mid\theta]\bigr]}_{\sigma^2_{\mathrm{aleatoric}}}.
    \label{eq:var_decomp}
\end{align}

\begin{itemize}[leftmargin=2em]
    \item $\sigma^2_{\mathrm{epistemic}}$: 데이터가 더 쌓이면 줄어드는 \textbf{감축 가능} 불확실성.
    \item $\sigma^2_{\mathrm{aleatoric}}$: 기상 자체의 본질적 변동성. 데이터로 줄일 수 없음.
\end{itemize}

운영자 대시보드에는 이 두 값을 분리 표시하여 ``모델이 학습을 덜 한 것인가''와 ``내일이 본질적으로 변동성 큰 날인가''를 구분할 수 있게 한다.

%=========================================================
\section{강화 1: AR(1) 관측 노이즈 --- 시간 상관관계 보존}
%=========================================================

\subsection{문제}

v1의 Beta 관측 모형은 각 시점을 독립으로 본다. 즉 $\tilde y_{i,t}$와 $\tilde y_{i,t+1}$의 noise는 독립. 그러나 실제 PV는 \textbf{구름이 서서히 지나가면서} 연속 시점에 걸쳐 함께 낮아진다. 독립 가정 하에서 시점별 분위수를 이어붙이면 ``램프 이벤트''가 희석되어, 운영상 가장 위험한 시나리오가 시나리오 번들에 포함되지 않는다.

\subsection{확장 모형}

잠재 노이즈 $\epsilon_{i,t}$에 AR(1) 구조를 둔다.

\begin{equation}
    \mathrm{logit}(\mu_{i,t}) = \eta_{i,t} + \epsilon_{i,t},\qquad
    \epsilon_{i,t} = \rho_r\,\epsilon_{i,t-1} + \nu_{i,t},\quad \nu_{i,t}\sim\mathcal{N}(0,\sigma^2_{\nu,r}).
    \label{eq:ar1_noise}
\end{equation}

$\rho_r\in(-1,1)$는 지역별 자기상관 계수이며, prior로 $\rho_r\sim\mathrm{Uniform}(-1,1)$ 또는 $\rho_r\sim\mathrm{Beta}(2,2)$의 $(-1,1)$ 재스케일링을 사용한다.

\subsection{효과}

\begin{itemize}[leftmargin=2em]
    \item \textbf{Full-trajectory sampling}: 알고리즘 \ref{alg:scenario}의 step 8에서 $\epsilon$ 계열을 AR(1)로 전진 생성하면, 한 시나리오 내부의 연속 시점이 자동으로 상관되어 램프가 살아난다.
    \item \textbf{예측 구간의 시간 누적}: horizon이 길수록 $\mathrm{Var}[\epsilon_{i,t+h}] = \sigma^2_{\nu,r}\cdot\frac{1-\rho_r^{2h}}{1-\rho_r^2}$로 커지므로 예측 구간이 시간에 따라 확장되는 현실적 형태를 띈다.
\end{itemize}

%=========================================================
\section{강화 2: Heteroscedastic Precision --- 날씨 의존 분산}
%=========================================================

\subsection{문제}

v1에서는 지역별 precision $\phi_r$가 상수이다. 그러나 \texttt{CLAUDE.md}에 정리된 현실은 다음과 같다.

\begin{table}[H]
\centering
\small
\caption{날씨별 PV 예측 오차 수준 (업계 관측치)}
\begin{tabular}{lll}
\toprule
날씨 & 오차 수준 & 빈도 \\
\midrule
맑은 날 & 2--5\% & 30--40\% \\
구름 조금 & 15--30\% & 30--40\% \\
흐린 날 & 20--40\% & 15--20\% \\
급변(Ramp) & 30--60\% & 5--10\% \\
\bottomrule
\end{tabular}
\end{table}

$\phi_r$가 상수라는 가정은 ``맑은 날 예측 구간''과 ``흐린 날 예측 구간''이 같다는 가정과 같아서, 운영자에게 잘못된 확신을 준다.

\subsection{확장 모형}

시점별 $\log \phi_{r,t}$를 날씨 변수의 선형함수로 둔다.
\begin{equation}
    \log\phi_{r,t} = \zeta_{r,0} + \zeta_{r,\mathrm{CLD}}\,\widetilde{\mathrm{CLD}}_{i,t} + \zeta_{r,\mathrm{CSI}}\,\widetilde{\mathrm{CSI}}_{i,t} + \zeta_{r,\mathrm{ramp}}\,|\Delta \widetilde{\mathrm{GHI}}_{i,t}|.
    \label{eq:hetero_phi}
\end{equation}

사전분포는
\begin{equation}
    \zeta_{r,0}\sim\mathcal{N}(\mu_\zeta,\sigma^2_\zeta),\quad \zeta_{r,\cdot}\sim\mathcal{N}(0,0.5^2).
\end{equation}

\subsection{효과}

\begin{itemize}[leftmargin=2em]
    \item \textbf{자동 캘리브레이션}: 구름 많은 시점은 $\phi$가 작아져 Beta 분포가 넓어지고, 모델이 ``자신 없음''을 분포 폭으로 표현한다.
    \item \textbf{운영상 의미}: 내일 전운량 예보가 높으면 $\phi_{r,t+h}$가 작아지고, 시나리오 번들의 분산이 자연스럽게 커지므로 하류 UC 모듈이 더 많은 예비력을 확보하도록 유도된다.
\end{itemize}

%=========================================================
\section{강화 3: Gaussian Process 공간 Prior --- 사이트간 공간 상관}
%=========================================================

\subsection{문제}

v1에서 사이트별 random intercept는 독립이다.
\begin{equation}
    u_i\sim\mathcal{N}(0,\sigma^2_u),\quad \mathrm{Cov}(u_i,u_j)=0.
\end{equation}

그러나 영흥수상태양광과 영흥육상태양광은 동일 기상장에 노출되므로 잔차가 매우 상관되어 있다. 독립 가정은 \textbf{포트폴리오 총 불확실성을 과소평가}한다. 이는 본 플랫폼의 핵심 결함 후보이다. 왜냐하면
\begin{equation}
    \mathrm{Var}\Bigl(\sum_i y_i\Bigr) = \sum_i \mathrm{Var}(y_i) + 2\sum_{i<j}\mathrm{Cov}(y_i,y_j),
\end{equation}
에서 공분산항을 0으로 가정하면 포트폴리오 분산이 작아 보이고, 그 결과 필요 예비력이 작아 보여 실시간에 부족 사태가 발생할 수 있기 때문이다.

\subsection{확장 모형}

사이트 좌표 $\bm\ell_i = (\mathrm{lat}_i,\mathrm{lon}_i)$를 사용해 GP prior를 둔다.
\begin{equation}
    \bm u = (u_1,\dots,u_N)^{\top} \sim \mathcal{N}(\bm 0, \bm K),\qquad
    K_{ij} = \sigma^2_u\,\exp\!\Bigl(-\frac{\lVert\bm\ell_i-\bm\ell_j\rVert^2}{2\ell_{\mathrm{GP}}^2}\Bigr).
    \label{eq:gp_prior}
\end{equation}

하이퍼파라미터 prior:
\begin{equation}
    \sigma_u\sim\mathrm{half\mbox{-}}t(3,0,1),\qquad \ell_{\mathrm{GP}}\sim\mathrm{InvGamma}(2,20\,\mathrm{km}).
\end{equation}

$\ell_{\mathrm{GP}}$ prior의 평균 20km는 ``구름 하나가 대체로 20km 반경에서 상관된다''는 기상학적 경험에 근거한다.

\subsection{계산 부담 완화}

사이트가 $N\sim 20$ 수준이면 full GP가 계산 가능하지만, 확장 시에는 다음을 고려한다.
\begin{itemize}[leftmargin=2em]
    \item \textbf{Kronecker 구조}: 공간 GP $\otimes$ 시간 AR(1) 분해
    \item \textbf{Predictive process / inducing points}: $M\ll N$개 inducing points로 rank 감축
    \item \textbf{Nearest Neighbor GP (NNGP)}: $O(N)$ 추론
\end{itemize}

남동발전 태양광 약 20개 사이트 규모에서는 full GP를 그대로 적용한다.

%=========================================================
\section{강화 4: Zero-Inflated Beta (ZIB) --- 일출/일몰 전이구간}
%=========================================================

\subsection{문제}

v1은 $d_t = \mathbb{I}(\mathrm{SZA}_t<90^\circ)$로 주간/야간을 결정적으로 분리한다. 문제는 일출 직후/일몰 직전 10--30분 구간이다. 이 구간에서는
\begin{itemize}[leftmargin=2em]
    \item 일부 사이트는 여전히 0 (그늘, 지형 차폐, 인버터 기동 대기)
    \item 일부 사이트는 매우 작은 양수 (0.1--5\% 수준)
\end{itemize}
가 섞여 있다. Beta 분포는 경계 $\{0,1\}$에 질량을 둘 수 없으므로, 이 구간에서 모델이 체계적 편향을 보인다.

\subsection{확장 모형: Zero-Inflated Beta (ZIB)}

\textbf{명칭 주의}: 이전 초안은 ``ZOIB (Zero-One-Inflated Beta)''로 표기하였으나, PV 정규화 출력 $\tilde y\in[0,1]$에서 1.0(설치용량 정확히 100\% 발전)은 사실상 발생하지 않으므로, 1에 대한 inflation은 불필요하다. 따라서 정확한 모형은 \textbf{Zero-Inflated Beta (ZIB)}이다.

\begin{equation}
    p(\tilde y_{i,t}\mid\cdot) = \pi_{0,i,t}\,\delta_0(\tilde y_{i,t}) + (1-\pi_{0,i,t})\,\mathrm{Beta}\bigl(\mu_{i,t}\phi_{r,t},\,(1-\mu_{i,t})\phi_{r,t}\bigr).
    \label{eq:zib}
\end{equation}

$\pi_{0,i,t}$는 ``정확히 0일 확률''이며 태양기하 변수의 함수로 모델링한다.
\begin{equation}
    \mathrm{logit}(\pi_{0,i,t}) = \omega_{r,0} + \omega_{r,\mathrm{SZA}}\,\mathrm{SZA}_t + \omega_{r,\mathrm{shade}}\,\mathbb{I}(\text{차폐 시간대}).
\end{equation}

전형적 주간 구간($\mathrm{SZA}\ll 90^\circ$)에서는 $\pi_0\approx 0$이 자동 학습되므로, 일반 Beta 회귀와 동일하게 작동한다. 경계 근처에서만 zero-inflation이 활성화된다.

\paragraph{향후 확장 (선택).} 만약 특정 사이트에서 inverter clipping으로 인해 $\tilde y \approx 1$이 빈번하게 관측된다면, 식 \eqref{eq:zib}을 다음과 같이 ZOIB로 확장할 수 있다.
\begin{equation*}
    p(\tilde y) = \pi_0\,\delta_0 + \pi_1\,\delta_1 + (1-\pi_0-\pi_1)\,\mathrm{Beta}(\cdot).
\end{equation*}
남동발전 사이트 데이터로 capacity factor 분포를 확인 후 결정한다.

\subsection{효과}

일출/일몰 1--2시간 구간의 예측 편향이 완화되고, 해당 시간대의 시나리오 분포가 ``많은 경우 0, 드물게 소량 발전''이라는 실제 패턴을 반영하게 된다.

%=========================================================
\section{강화된 1단계 통합 식}
%=========================================================

v2의 완전한 1단계 모형은 다음과 같이 종합된다. \textbf{$\epsilon$은 $\eta$ 안에 포함시키지 않고 link 단계에서 가산}하는 형태로 정리하여 v1과의 비교를 명확히 한다.

\begin{align}
    \eta_{i,t} &= \underbrace{\alpha_0 + a_{r(i)} + u_i}_{\text{평균 + 공간 (GP prior on } \bm u)}
    + \underbrace{(\beta_{\mathrm{GHI}} + b_{r(i),\mathrm{GHI}})\widetilde{\mathrm{GHI}}_{i,t} + (\beta_T + b_{r(i),T})\widetilde T_{i,t} + (\beta_{\mathrm{CLD}}+b_{r(i),\mathrm{CLD}})\widetilde{\mathrm{CLD}}_{i,t} + \cdots}_{\text{기상 random slope}} \nonumber\\
    &\quad + \underbrace{\sum_{k=1}^{K_h}\bigl[\gamma_k^{(s)}\sin + \gamma_k^{(c)}\cos\bigr] + \sum_{m=1}^{K_d}\bigl[\delta_m^{(s)}\sin + \delta_m^{(c)}\cos\bigr]}_{\text{전역 Fourier (일/연)}} \nonumber\\
    &\quad + \underbrace{\sum_{k=1}^{K_h}\bigl[\tilde\gamma_{r,k}^{(s)}\sin + \tilde\gamma_{r,k}^{(c)}\cos\bigr] + \sum_{m=1}^{K_d}\bigl[\tilde\delta_{r,m}^{(s)}\sin + \tilde\delta_{r,m}^{(c)}\cos\bigr]}_{\text{지역별 Fourier 편차 (shrinkage)}}
    \label{eq:bayes_stage1_v2}
\end{align}

\begin{align}
    \mathrm{logit}(\mu_{i,t}) &= \eta_{i,t} + \epsilon_{i,t}, &
    \epsilon_{i,t} &= \rho_r\epsilon_{i,t-1} + \nu_{i,t},\quad \nu_{i,t}\sim\mathcal{N}(0,\sigma^2_{\nu,r}) \quad\text{(식 \eqref{eq:ar1_noise})}, \\
    \log\phi_{r,t} &= \zeta_{r,0} + \zeta_{r,\mathrm{CLD}}\widetilde{\mathrm{CLD}}_{i,t} + \zeta_{r,\mathrm{CSI}}\widetilde{\mathrm{CSI}}_{i,t} + \zeta_{r,\mathrm{ramp}}|\Delta\widetilde{\mathrm{GHI}}_{i,t}|, & & \text{(식 \eqref{eq:hetero_phi})} \\
    \mathrm{logit}(\pi_{0,i,t}) &= \omega_{r,0} + \omega_{r,\mathrm{SZA}}\mathrm{SZA}_t + \omega_{r,\mathrm{shade}}\mathbb{I}_{\mathrm{shade}}, & & \\
    \tilde y_{i,t} &\sim \mathrm{ZIB}\bigl(\pi_{0,i,t},\,\mu_{i,t},\,\phi_{r,t}\bigr) & & \text{(식 \eqref{eq:zib})}, \\
    \bm u &\sim \mathcal{N}(\bm 0, \bm K), & K_{ij} &= \sigma^2_u\exp\!\bigl(-\lVert\bm\ell_i-\bm\ell_j\rVert^2/(2\ell_{\mathrm{GP}}^2)\bigr) \quad\text{(식 \eqref{eq:gp_prior})}.
\end{align}

즉, v1의 식 \texttt{(eq:bayes\_stage1)}에 \textbf{AR(1) 잠재 노이즈, 이분산 $\phi$, GP 공간 prior, ZIB 관측분포}가 추가된 형태이다.

\paragraph{식 검증 — v1 특수 케이스 회수.}
$\rho_r=0,\ \sigma^2_{\nu,r}=0,\ \zeta_{r,\cdot}=0,\ \omega_{r,\cdot}\to-\infty,\ \ell_{\mathrm{GP}}\to\infty$로 두면 본 식 (\ref{eq:bayes_stage1_v2})이 v1의 식 (\texttt{eq:bayes\_stage1})로 환원된다. 즉 v2 모형은 v1을 \textbf{진짜 특수 케이스로 포함하는 상위집합}임이 식 수준에서 확인된다.

%=========================================================
\section{1단계 출력물 재정의 (강화판)}
%=========================================================

v2의 1단계는 다음 다섯 종류의 산출물을 생성한다.

\begin{table}[H]
\centering
\small
\caption{1단계 출력물 분류}
\begin{tabular}{p{2.5cm}p{4.5cm}p{7cm}}
\toprule
범주 & 내용 & 하류 사용처 \\
\midrule
A. 점추정 & $\hat\mu^{\mathrm{Bayes}}_{i,t+h}$, 잔차 $e_{i,t}$ & 2단계 FiLM-GRU 입력, 대시보드 중앙값 \\
B. 시나리오 번들 & $\{\tilde y^{(s)}\}_{s=1}^{S}$, shape $(S,H,N)$ & \textbf{Stochastic UC 직접 입력} \\
C. 분위수 요약 & $P_5, P_{10}, P_{50}, P_{90}, P_{95}$; $\mathrm{CVaR}_5, \mathrm{CVaR}_{95}$ & 운영자 대시보드, KPX 신고 밴드 \\
D. 불확실성 분해 & $\sigma_{\mathrm{epistemic}}, \sigma_{\mathrm{aleatoric}}(t)$ & ``학습 부족'' vs ``본질적 변동'' 구분 \\
E. Regional signature & 식 \eqref{eq:regional_signature_v2} & 2단계 FiLM conditioning \\
\bottomrule
\end{tabular}
\end{table}

\subsection{Regional signature 확장 (E)}

v1의 regional signature에 v2의 새 파라미터를 추가한다.

\begin{align}
    \bm z_r^{\mathrm{stat,v2}} = \Big[&
    \underbrace{\mathbb{E}[a_r],\mathbb{E}[b_{r,\mathrm{GHI}}],\mathbb{E}[b_{r,T}],\mathbb{E}[b_{r,\mathrm{CLD}}]}_{\text{v1: 지역 평균/민감도}},\;
    \underbrace{\mathbb{E}[\tilde\gamma_{r,\cdot}],\mathbb{E}[\tilde\delta_{r,\cdot}]}_{\text{v1: 시간/계절 구조}}, \nonumber\\
    &\underbrace{\mathbb{E}[\rho_r],\mathbb{E}[\sigma_{\nu,r}]}_{\text{v2: 시간 상관}},\;
    \underbrace{\mathbb{E}[\zeta_{r,\cdot}]}_{\text{v2: 이분산}}, \nonumber\\
    &\underbrace{\mathrm{sd}(a_r),\mathrm{sd}(b_{r,\cdot}),\mathrm{sd}(\rho_r),\dots}_{\text{posterior uncertainty}}
    \Big].
    \label{eq:regional_signature_v2}
\end{align}

즉, 2단계 FiLM은 ``이 지역이 얼마나 시간 상관이 강한지''와 ``이 지역이 얼마나 날씨 의존적 이분산인지''까지 조건 신호로 받는다.

%=========================================================
\section{2단계 확장: Probabilistic Residual Head}
%=========================================================

\subsection{결정론적 head의 한계}

v1의 2단계는 잔차를 점예측한다.
\begin{equation}
    \hat e_{i,t+h} = f_{\mathrm{out}}^{(h)}(\bm u_{i,t}).
\end{equation}
이 경우 최종 예측 $\hat{\tilde y}_{i,t+h} = \hat\mu^{\mathrm{Bayes}}_{i,t+h} + \hat e_{i,t+h}$도 점예측이다. 시나리오 관점에서는 2단계가 1단계의 불확실성을 ``부분적으로만'' 전파하게 된다.

\subsection{Probabilistic head}

2단계 출력을 평균과 분산 두 채널로 둔다.
\begin{equation}
    (\hat e_{i,t+h},\,\hat\sigma_{e,i,t+h}^2) = f_{\mathrm{out}}^{(h)}(\bm u_{i,t}).
\end{equation}

학습 손실을 Gaussian NLL로 한다.
\begin{equation}
    \mathcal{L}_{\mathrm{nll}} = \sum_h w_h\,\Bigl[\frac{(e_{i,t+h}-\hat e_{i,t+h})^2}{2\hat\sigma_{e,i,t+h}^2} + \frac{1}{2}\log\hat\sigma_{e,i,t+h}^2\Bigr].
\end{equation}

최종 시나리오는 1단계 시나리오 $\tilde y^{(s)}$에 2단계 잔차 샘플을 추가로 합성하여 생성한다.
\begin{equation}
    \hat y^{(s,k)}_{i,t+h} = C_i\Bigl[\tilde y^{(s)}_{i,t+h} + \hat e_{i,t+h} + \hat\sigma_{e,i,t+h}\cdot\xi^{(k)}\Bigr],\quad \xi^{(k)}\sim\mathcal{N}(0,1).
\end{equation}

$S$개의 1단계 시나리오 $\times$ $K$개의 2단계 잔차 perturbation으로 총 $S\cdot K$개 최종 시나리오가 얻어진다. 전형적 설정은 $S=500, K=2$이다.

%=========================================================
\section{하류 연계: Stochastic UC와 Chance-Constrained UC}
%=========================================================

\subsection{점예측 UC (기존 단순안)}

\begin{align}
    \min_{P_g}\quad & \sum_g c_g^{\mathrm{var}} P_{g,t} \\
    \text{s.t.}\quad & \sum_g P_{g,t} + \hat{\mathrm{PV}}_t = D_t, \\
    & P_g^{\min}\le P_{g,t}\le P_g^{\max}, \quad \text{ramp 제약}, \dots
\end{align}

$\hat{\mathrm{PV}}$가 틀리면 실시간에 비싼 LNG 급기동이 필요하고, 이것이 남동발전 백업 비용의 핵심 낭비 요인이다.

\subsection{Stochastic UC (v2 제안)}

시나리오 번들 $\{\mathrm{PV}^{(s)}_t\}_{s=1}^{S}$를 직접 입력으로 받는다.

\begin{align}
    \min_{u_g, P_g^{(s)}}\quad & \underbrace{\sum_g c_g^{\mathrm{SU}} u_{g,t}}_{\text{기동비용 (here-and-now)}} + \frac{1}{S}\sum_{s=1}^{S}\underbrace{\sum_g c_g^{\mathrm{var}} P_{g,t}^{(s)}}_{\text{시나리오별 가변비 (wait-and-see)}} \\
    \text{s.t.}\quad & \sum_g P_{g,t}^{(s)} + \mathrm{PV}^{(s)}_t = D_t,\quad \forall s, \\
    & P_g^{\min} u_{g,t}\le P_{g,t}^{(s)}\le P_g^{\max} u_{g,t},\quad \forall s, \\
    & \text{ramp, 최소기동시간, 최소정지시간 제약}.
\end{align}

기동 결정 $u_g$는 시나리오에 무관하게 공통(``here-and-now''), 각 시나리오 하의 발전량 $P_g^{(s)}$는 시나리오별로 달라짐(``wait-and-see'').

\subsection{Chance-Constrained UC (해석 친화형)}

운영자에게 더 직관적인 표현:
\begin{equation}
    \mathbb{P}\bigl(\mathrm{PV}_t < \underline{\mathrm{PV}}_t\bigr) \le \varepsilon\quad\Longleftrightarrow\quad \text{화력 예비력} \ge \hat y^{P_{50}}_t - \hat y^{P_{\varepsilon\cdot 100}}_t.
\end{equation}

$\varepsilon=0.05$라 두면, ``P5 시나리오까지 커버할 만큼의 석탄 가용용량을 신고''하는 것이 곧 KPX 10시 신고값이 된다. v2의 시나리오 출력은 이 값을 \textbf{직접} 산출한다.

\begin{keybox}[핵심 연계]
v2의 PV 예측 모듈이 산출하는 분위수 $\hat y^{P_5}$가 곧 \textit{plan v6}의 ``KPX 가용용량 신고 최적화''의 정량적 답변이다. 즉, PV 불확실성이 화력 예비력 산정에 직접 변환된다.
\end{keybox}

%=========================================================
\section{평가 프레임워크 확장}
%=========================================================

v1의 평가지표(nRMSE, nMAE)는 시나리오 품질을 평가할 수 없다. v2는 확률예측 표준 지표를 추가한다.

\subsection{점예측 평가 (v1 유지)}
\begin{itemize}[leftmargin=2em]
    \item nRMSE, nMAE: 중앙값 예측의 정확도
    \item Skill Score vs Persistence
\end{itemize}

\subsection{확률예측 평가 (v2 신규)}

\textbf{CRPS (Continuous Ranked Probability Score)}:
\begin{equation}
    \mathrm{CRPS}(F,y) = \int_{-\infty}^{\infty}\bigl(F(z) - \mathbb{I}(y\le z)\bigr)^2 dz.
\end{equation}
샘플 근사: $\mathrm{CRPS}(\{y^{(s)}\},y) = \frac{1}{S}\sum_s|y^{(s)}-y| - \frac{1}{2S^2}\sum_{s,s'}|y^{(s)}-y^{(s')}|$.

\textbf{Coverage}: $P_{10}$--$P_{90}$ 구간이 실제로 80\%를 포함하는 비율. 이상적으로 0.80.

\textbf{Winkler Score}: interval forecast 품질. 좁고 정확할수록 낮음.

\textbf{Pinball Loss}: 분위수별 비대칭 손실. 여러 $\tau$에 대해 보고.

\subsection{운영 지표 (v2 신규, 가장 중요)}

\textbf{Ramp Event Hit Rate}: 실제 발생한 급변(예: 1시간 내 20\% 이상 하락)이 $P_5$--$P_{50}$ 밴드 안에 있었는지 비율. 목표 $\ge 0.95$.

\textbf{Downstream UC Cost}: 시나리오 번들을 Stochastic UC에 넣고 나온 total operating cost 대 점예측 UC 대비 절감율. \textbf{본 플랫폼의 최종 KPI}.

\textbf{Reserve Shortfall Frequency}: P5 기반 예비력 신고 시 실시간에 발전 부족이 발생한 빈도. 이상적 $\le 5\%$.

\subsection{Ablation 비교군 (v2 확장)}

\begin{table}[H]
\centering
\small
\caption{v2 ablation 구성 (v1 B1--B8을 확장)}
\begin{tabular}{ll}
\toprule
ID & 설명 \\
\midrule
A1 & v1 모형 (basic Bayesian + GRU + FiLM + signature) \\
A2 & A1 + AR(1) noise \\
A3 & A1 + heteroscedastic $\phi$ \\
A4 & A1 + GP spatial prior \\
A5 & A1 + ZIB \\
A6 & A1 + 모든 v2 강화 (full v2) \\
A7 & A6 + 2단계 probabilistic head \\
\midrule
\multicolumn{2}{c}{\textbf{평가: 점예측 + 확률예측 + 운영 지표 모두}} \\
\bottomrule
\end{tabular}
\end{table}

%=========================================================
\section{구현 로드맵}
%=========================================================

\subsection{Phase 구조 재편}

\begin{longtable}{p{1.2cm}p{5cm}p{8cm}}
\toprule
Phase & 제목 & 주요 산출물 \\
\midrule
\endhead
0 & 데이터 수집 및 EDA & 남동발전 크롤링, 기상청 API, 사이트 메타 (v1과 동일) \\
1 & v1 기준 1단계 PyMC 구현 & Basic Bayesian 학습, posterior 진단, 시나리오 1,000개 생성 \\
2 & 강화 1--4 순차 적용 & A2--A5 각각 학습, CRPS/coverage 측정, 효과 확인 \\
3 & 통합 v2 1단계 (A6) & 완전 v2 1단계 학습 및 A1 대비 평가 \\
4 & 2단계 FiLM-GRU + probabilistic head & A7 학습, 최종 시나리오 생성 파이프라인 완성 \\
5 & Stochastic UC 연계 & \texttt{plan/fuel/}로 이관. PV 시나리오 입력 인터페이스 정의 \\
6 & 실시간 운영 모드 & 초단기예보/위성 기반 실시간 보정, 1단계 고정 + 2단계 실시간 업데이트 \\
\bottomrule
\end{longtable}

\subsection{우선순위 근거}

강화 적용 순서는 \textbf{운영 임팩트 기준}으로 정한다.

\begin{enumerate}[leftmargin=2em]
    \item \textbf{강화 3 (GP 공간 prior)}: 포트폴리오 총 분산 추정을 즉각 현실화. 현재 구조는 총 불확실성을 과소평가하고 있을 가능성이 높음.
    \item \textbf{강화 2 (heteroscedastic $\phi$)}: 맑은 날/흐린 날 신뢰도 차별화. 운영자 대시보드에 바로 반영됨.
    \item \textbf{강화 1 (AR(1) noise)}: 램프 이벤트 시나리오 품질. 운영상 가장 위험한 케이스를 잡는 데 결정적.
    \item \textbf{강화 4 (ZIB)}: 일출/일몰 구간 편향. 수익 임팩트는 작지만 모델 정합성을 위해 필요.
\end{enumerate}

%=========================================================
\section{기술 스택 보강}
%=========================================================

\begin{table}[H]
\centering
\small
\caption{v2 기술 스택}
\begin{tabular}{lll}
\toprule
용도 & 도구 & 비고 \\
\midrule
베이지안 1단계 & \textbf{NumPyro} or PyMC 5 & NumPyro는 JAX 기반 고속. GP + AR(1) 결합 시 유리 \\
GP 근사 & GPyTorch or NumPyro 내장 & NNGP 확장 가능성 \\
딥러닝 2단계 & PyTorch & v1과 동일 \\
Probabilistic head & torch.distributions & Gaussian NLL 구현 \\
확률예측 평가 & \textbf{properscoring}, scoringrules & CRPS, Brier 등 \\
시나리오 기반 UC & \textbf{Pyomo + Gurobi/CBC} & \texttt{plan/fuel/}에서 구현 \\
실험관리 & MLflow & v1과 동일 \\
\bottomrule
\end{tabular}
\end{table}

\texttt{PyMC $\to$ NumPyro} 전환을 고려하는 이유는 GP prior와 AR(1) 잠재변수가 결합되면 MCMC가 느려지기 때문이다. NumPyro의 NUTS는 JAX jit 컴파일로 10배 이상 빠른 경우가 많다.

%=========================================================
\section{v1에서 v2로 마이그레이션 체크리스트}
%=========================================================

기존 v1 작업과의 호환성 유지를 위한 체크리스트.

\begin{itemize}[leftmargin=2em]
    \item[$\square$] v1의 Phase 0 (데이터 수집) 완전히 재사용
    \item[$\square$] v1의 Phase 1-1 (데이터 준비) 전처리 로직 재사용, ZIB용 경계 변환은 추가
    \item[$\square$] v1의 지역 클러스터링 결과 재사용, GP prior는 원시 위경도 사용
    \item[$\square$] v1의 PyMC 코드를 NumPyro로 포팅 (선택)
    \item[$\square$] v1의 regional signature에 $\rho_r, \sigma_{\nu,r}, \zeta_{r,\cdot}$ 추가
    \item[$\square$] v1의 2단계 FiLM-GRU에 probabilistic head 추가 (출력 차원 $H \to 2H$)
    \item[$\square$] v1의 평가 스크립트에 CRPS, coverage, Winkler, pinball 추가
    \item[$\square$] \texttt{plan/fuel/}와의 인터페이스 문서화 (시나리오 텐서 shape, 단위, 시간대)
\end{itemize}

%=========================================================
\section{우려사항과 대응}
%=========================================================

\begin{table}[H]
\centering
\small
\caption{v2 구현 시 우려사항 및 대응}
\begin{tabular}{p{4cm}p{9.5cm}}
\toprule
우려 & 대응 \\
\midrule
MCMC 계산 부담 증가 & NumPyro + GPU, variational inference (SVI) 보조, 사이트 수 제한된 범위에서 full GP 유지 \\
GP 공분산 행렬 특이성 & nugget term $\tau^2 I$ 추가 ($\tau^2\sim 10^{-6}$) \\
ZIB의 $\pi_0$ 식별성 & SZA 경계 구간에서만 학습되도록 데이터 가중치 조정 \\
시나리오 개수 $S$ 선택 & CRPS 수렴까지 $S$ 증가. 통상 $S=500$--1000에서 충분 \\
2단계 probabilistic head 과적합 & monte carlo dropout 또는 deep ensemble 보강 \\
운영자 혼란 (``왜 1000개?'') & 대시보드는 분위수 밴드와 중앙값만 노출, 시나리오는 UC에만 공급 \\
\bottomrule
\end{tabular}
\end{table}

%=========================================================
\section{요약}
%=========================================================

\begin{keybox}[v2의 한 문장 요약]
\textit{v2는 PV 예측 모듈을 ``점예측기''가 아니라 ``불확실성 생성기''로 재정의하여, 베이지안 posterior에서 직접 추출한 1,000개의 포트폴리오 시계열 시나리오를 하류 Stochastic UC 모듈에 공급함으로써, 태양광 변동성을 화력 예비력 결정에 정량적으로 연결한다.}
\end{keybox}

\subsection{기여 문장}
\begin{enumerate}[label=\textbf{C\arabic*.}, leftmargin=2.2em]
    \item 베이지안 계층모형의 출력을 점예측이 아닌 \emph{시나리오 번들}로 재정의하여, PV 예측과 화력 최적화를 확률 기반으로 직렬 연결하는 프레임워크를 제시한다.
    \item 기존 Beta 회귀 관측모형에 AR(1) 잠재 노이즈, heteroscedastic precision, Gaussian Process 공간 prior, Zero-Inflated Beta (ZIB)를 결합하여, 시간/공간/날씨 의존적 불확실성을 명시적으로 모델링한다.
    \item 2단계 FiLM-GRU에 probabilistic residual head를 도입하여 1단계의 불확실성이 2단계를 거친 후에도 보존되도록 한다.
    \item 최종 산출물로 분위수 요약, 불확실성 분해, 원시 시나리오 번들의 삼중 출력 구조를 정의하고, 각 하류 소비자(대시보드, KPX 신고, Stochastic UC)별 사용처를 명시한다.
    \item 평가 프레임워크를 점예측 지표에서 확률예측 지표(CRPS, coverage, Winkler, pinball) 및 운영 지표(ramp hit rate, downstream UC cost)로 확장한다.
\end{enumerate}

\subsection{다음 단계}

\begin{enumerate}[leftmargin=2em]
    \item \texttt{plan/fuel/} 작성: PV 시나리오 번들을 소비하는 Stochastic UC 설계서
    \item \texttt{plan/pv/v2} 기반으로 Phase 1--3 구현 착수 (데이터 확보 전제)
    \item 강화 3 (GP spatial prior) 우선 프로토타이핑으로 포트폴리오 총 불확실성 크기 실측
\end{enumerate}

\end{document}
